% from ega4-8-fr.tex
\oldpage[IV-3]{38}
\item[(vii)] radiciel~;
\item[(viii)] affine~;
\item[(ix)] quasi-affine~;
\item[(x)] fini~;
\item[(xi)] quasi-fini~;
\item[(xii)] propre.
\end{enumerate}
Alors, si $\mathbf P$ désigne une des propriétés précédentes, pour que $f$
possède la propriété $\mathbf P$, il faut et il suffit qu'il existe
$\lambda\geq\alpha$ tel que $f_\lambda$ possède la propriété $\mathbf P$
(auquel cas $f_\mu$ la possède aussi pour $\mu\geq\lambda$).

Si $S_0$ est en outre supposé quasi-séparé, la même conclusion est valable
lorsque $\mathbf P$ est la propriété d'être~:
\begin{enumerate}
\item[(xiii)] projectif~;
\item[(xiv)] quasi-projectif.
\end{enumerate}
\end{theorem}

\begin{proof}
Le cas où $\mathbf P$ est l'une des propriétés (i) ou (v) n'a été inséré dans
l'énoncé que pour mémoire~; dans ces cas, le théorème découle de ce qui a été
démontré respectivement dans (8.8.2.4) et (8.10.4). D'ailleurs, compte tenu de
(I, 5.4.1 et 5.3.4), (v) résulte aussi de (iv). Le cas (i \emph{bis}) se
déduit aussitôt de (i), en utilisant (I, 5.3.8) et en notant (comme on l'a
déjà utilisé dans (8.10.4)) que la diagonale $\Delta_f$ se déduit de
$\Delta_{f_\lambda}$ par le changement de base $S\to S_\lambda$.

On notera d'autre part que (vi), (vii) et (xi) sont en fait des conditions
sur les fibres $f^{-1}(y)$ des morphismes envisagés, compte tenu de la
transitivité des fibres par changement de base (I, 3.6.4)~: la condition (vi)
signifie en effet que toutes les fibres doivent être non vides, la condition
(vii) qu'elles doivent être radicielles (I, 3.5.8) et la condition (xi)
qu'elles doivent être finies (III, 6.2.2 et 6.2.3 et I, 6.4.4, compte tenu de
ce que $f$ et les $f_\lambda$ sont des morphismes de type fini par (I, 5.4,
(v)). Le théorème dans ces trois cas sera donc encore conséquence d'un
résultat général sur ce type de propriétés ne concernant que les fibres, qui
sera établi dans (9.3.3)~; nous repoussons donc jusqu'à ce moment-là la
démonstration du théorème dans le cas (xi) (bien entendu, le lecteur pourra
vérifier que, sauf aux n\textsuperscript{os}~8.11 et 8.12, nous ne ferons pas
usage du théorème dans ce cas jusqu'à (9.3.3) et que (8.11) et (8.12) ne
seront pas utilisés avant (9.3.3)).

Pour les cas qui restent à démontrer, on peut se borner à prouver que la
condition de l'énoncé est nécessaire, toutes les propriétés $\mathbf P$
envisagées étant invariantes par changement de base (voir chap. I\textsuperscript{er}
et II dans les numéros concernant chacune de ces propriétés). On peut en
outre supposer que $S_\alpha=S_0$ et que $Y_\alpha=S_\alpha$, donc
$Y_\lambda=S_\lambda$ pour tout $\lambda\geq\alpha$. Enfin, les propriétés
(i) à (xii) sont locales sur $S_0$, donc, comme $S_0$ est réunion finie
d'ouverts affines et $L$ filtrant, on peut se borner pour les démontrer au cas
où $S_0=\operatorname{Spec}(A_0)$ est affine (donc quasi-séparé). On désignera
par $A_\lambda$ (resp. $A$) l'anneau de $S_\lambda$ (resp. $S$).

\emph{Cas (ii), (iii), (iv)}~: Supposons que $f$ soit une immersion (resp.
une immersion ouverte, resp. une immersion fermée), et soit $X'$ le
sous-préschéma (resp. induit sur un ouvert, resp. fermé) de $S$ associé à $f$,
qui est donc un $S$-préschéma de présentation finie.
